\section{Differential Forms and the Exterior Derivative}
    \label{sec:forms-derivative}

    \subsection{Differential 1-Forms}
    In the previous section, we introduced the cotangent space $T_p^*M$ at a
    point $p$. Each element of this space is a linear map that takes a tangent
    velocity at $p$ and returns a real number.\footnote{%
        For an alternative, informal introduction to differential forms from
        the viewpoint of integration, see Tao \cite{tao2007differential}.
        The PDF link is included in the bibliography.}

    A \textbf{differential $1$-form} $\omega$ on $M$ assigns a cotangent
    vector $\omega_p \in T_p^*M$ to every point $p$, with this assignment
    varying smoothly from point to point.\footnote{%
        We leave the precise meaning of ``varying smoothly'' as a technical
        detail.} Thus, at each point, we have a
    linear measurement
    \begin{equation}
        \omega_p \colon T_pM \to \mathbb{R}.
    \end{equation}

    We already have an example: the differential $df_p$ at each point.
    Assigning these differentials across the manifold gives the $1$-form $df$.
    \textbf{Thus, $df_p$ is a measurement at one point, while $df$ provides
    these measurements throughout the manifold.}

    Returning to temperature on the sphere, $df_p([\gamma])$ tells us how
    quickly the temperature changes as a particle passes through $p$ along
    $\gamma$. The $1$-form $df$ provides this temperature-change measurement
    everywhere on the sphere.

    \paragraph{Vector Fields and $1$-Forms in Euclidean Space.}
    Recall that at each $p \in \mathbb{R}^n$, the cotangent vector $df_p$ corresponds
    to the vector $\nabla f(p)$ under the Euclidean dot-product identification:
    $df_p(v) = \nabla f(p) \cdot v$. Applying this at every point, the
    $1$-form $df$ corresponds to the gradient vector field $\nabla f$.
    More generally, any smooth vector field $F$ on $\mathbb{R}^n$ defines an
    associated $1$-form $\omega_F$ by
    \begin{equation}
        (\omega_F)_p(v) = F(p) \cdot v, \qquad v \in T_p\mathbb{R}^n = \mathbb{R}^n.
    \end{equation}
    Thus $df = \omega_{\nabla f}$. If $F$ is a force field and $v$ is a unit
    direction, $(\omega_F)_p(v)$ is the force component in that direction.

    \subsection{Geometric Motivation for Higher Forms}
    A $1$-form can measure change along a single tangent direction, as $df$
    does. To measure quantities associated with an area, as in the study of
    curl, we use higher forms. A $2$-form takes two tangent vectors
    $v,w \in T_pM$ and gives a signed measurement of the parallelogram they span. Swapping the vectors changes the sign of this measurement,
    just as swapping the vectors in a cross product in $\mathbb{R}^3$
    changes its sign:
    \[
        \omega_p(v,w) = -\omega_p(w,v).
    \]

    \subsection{Multilinear and Alternating Maps}
    To describe these measurements, we use multilinear and alternating maps
    on a vector space $V$, following Munkres \cite[\S26, p.~220;
    \S27, pp.~229--230]{munkres1991analysis}.

    \begin{definition}[Multilinear and Alternating Maps %
        \cite{munkres1991analysis}]
        Let $V$ be a real vector space. A map
        \[
            T \colon \underbrace{V \times \dots \times V}_{k \text{ copies}}
            \to \mathbb{R}
        \]
        is \textbf{multilinear of degree $k$} (or a \textbf{$k$-linear map}) if
        it is linear in each argument when all other $k-1$ arguments are held
        fixed.

        A multilinear map $T$ of degree $k$ is \textbf{alternating} if swapping
        any two distinct arguments reverses its sign: for any $1 \le i < j \le
        k$,
        \begin{equation}
            T(v_1, \dots, v_i, \dots, v_j, \dots, v_k)
            = -T(v_1, \dots, v_j, \dots, v_i, \dots, v_k).
        \end{equation}
        Consequently, if any two input vectors are identical ($v_i = v_j$ for
        $i \neq j$), an alternating map necessarily evaluates to zero:
        $T(\dots, v, \dots, v, \dots) = 0$.
    \end{definition}

    \subsection{Higher Differential Forms: At a Point and Globally}
    We can now define differential forms of arbitrary degree $k$ on a
    manifold, following Lee \cite[``Differential Forms on Manifolds,''
    pp.~359--360]{lee2013smooth}:
    \begin{enumerate}
        \item \textbf{At a single point $p \in M$}:\\
            A \textbf{$k$-form at $p$} is an alternating multilinear map on
            the tangent space $T_pM$:
            \[
                \omega_p \colon
                \underbrace{T_pM \times \dots \times T_pM}_{k \text{ copies}}
                \to \mathbb{R}.
            \]
        \item \textbf{On the manifold $M$ as a whole}:\\
            A \textbf{differential $k$-form} $\omega$ on $M$ is a smooth
            choice of an alternating multilinear map at each point
            \cite[p.~360]{lee2013smooth}.\footnote{%
                As with $1$-forms, ``smooth choice'' is an informal
                description; we do not formalize this notion in the talk.}
    \end{enumerate}

    We denote the space of smooth $k$-forms on $M$ by $\Omega^k(M)$, with
    $\Omega^0(M) = C^\infty(M)$ (smooth scalar functions).
    This is a real vector space: we add forms and multiply them by real
    scalars point by point. For $\omega,\omega' \in \Omega^k(M)$,
    $a \in \mathbb{R}$, and $v_1,\dots,v_k \in T_pM$,
    \[
        \begin{aligned}
            (\omega_p+\omega'_p)(v_1,\dots,v_k)
                &= \omega_p(v_1,\dots,v_k)+\omega'_p(v_1,\dots,v_k), \\
            (\omega+\omega')_p &= \omega_p+\omega'_p, \\
            (a\omega)_p &= a\omega_p.
        \end{aligned}
    \]
    Sums and scalar multiples remain smooth and alternating.

    \subsection{The Exterior Derivative}
    The differential of a smooth function turns a $0$-form into a $1$-form.
    This operation extends to differential forms of every degree through an
    operator called the \textbf{exterior derivative}:
    \begin{equation}
        d \colon \Omega^k(M) \to \Omega^{k+1}(M).
    \end{equation}
    We omit its construction (see \cite[Theorem~14.24, pp.~365--366]{lee2013smooth})
    and retain three properties needed for this talk:
    \begin{enumerate}
        \item \textbf{Linearity}:\\
            For $a,b \in \mathbb{R}$ and $\alpha,\beta \in \Omega^k(M)$,
            \[
                d(a\alpha + b\beta) = a\,d\alpha + b\,d\beta.
            \]
        \item \textbf{On functions}:\\
            For a smooth function $f$, $df$ is the differential introduced in
            Section~\ref{sec:tangent-cotangent}.
        \item \textbf{Differentiating twice gives zero}:\\
            For every differential form $\omega$,
            \[
                d(d\omega) = 0, \qquad \text{that is, } d^2 = 0.
            \]
    \end{enumerate}

    In $\mathbb{R}^3$, the successive exterior derivatives correspond to the
    familiar vector-calculus operators:
    \[
        \begin{aligned}
            d &\colon C^\infty(\mathbb{R}^3) \to \Omega^1(\mathbb{R}^3)
                &&\text{gradient}, \\
            d &\colon \Omega^1(\mathbb{R}^3) \to \Omega^2(\mathbb{R}^3)
                &&\text{curl}, \\
            d &\colon \Omega^2(\mathbb{R}^3) \to \Omega^3(\mathbb{R}^3)
                &&\text{divergence}.
        \end{aligned}
    \]
    The first map sends $f$ to $df$, which we already identified with the
    gradient in Section~\ref{sec:tangent-cotangent}. These are the familiar
    vector-calculus operators expressed in the language of differential forms
    on the manifold $\mathbb{R}^3$.\footnote{%
        For the precise identifications, see Lee
        \cite[Equation~(14.27), p.~368; Exercise~14.28]{lee2013smooth}.
        For explicit formulas, see also Bouchard
        \cite[\S4.4.1, Definitions~4.4.1--4.4.3 and
        Table~4.4.8]{bouchard_exterior_vector}.}

    There are no nonzero forms of degree greater than $3$ on $\mathbb{R}^3$;
    in particular, $\Omega^4(\mathbb{R}^3)=\{0\}$.\footnote{%
        Any four vectors in a three-dimensional tangent space are linearly
        dependent, so every alternating $4$-linear form vanishes. The same
        argument applies to all higher degrees.}
    Thus, the next map $d\colon\Omega^3(\mathbb{R}^3)\to\Omega^4(\mathbb{R}^3)$
    is zero.

    The familiar identities
    \[
        \operatorname{curl}(\nabla f) = 0,
        \qquad
        \operatorname{div}(\operatorname{curl} F) = 0
    \]
    are instances of $d^2 = 0$. This identity is the starting point for de Rham
    cohomology.
